\documentclass[11pt,a4paper]{article}
\usepackage[spanish,es-noquoting]{babel}
\usepackage[top=2.2cm, bottom=2.2cm, left=2.2cm, right=2.2cm]{geometry}
\usepackage{fontspec}
\usepackage{amsmath,amssymb}
\usepackage{enumitem}
\usepackage{microtype}

\setmainfont{Calibri}
\linespread{1.15}

\begin{document}
\pagestyle{empty}

\noindent\textbf{1.} Se desea fabricar un depósito cilíndrico cerrado con una capacidad de $V = 16\pi\text{ m}^3$. Materiales: tapa superior a $3\text{ €/m}^2$, base inferior a $5\text{ €/m}^2$ y pared lateral a $4\text{ €/m}^2$.
\begin{enumerate}[label=\textbf{\alph*)}, leftmargin=2em, itemsep=6pt]
    \item \textbf{Función del coste total $C(r)$:}
    \begin{itemize}[leftmargin=1.5em, itemsep=2pt]
        \item Ligadura del volumen: $V = \pi r^2 h = 16\pi \implies h = \frac{16}{r^2}$.
        \item Coste total:
        \[
        C(r) = 3(\pi r^2) + 5(\pi r^2) + 4(2\pi r h) = 8\pi r^2 + 8\pi r \left(\frac{16}{r^2}\right) = \mathbf{8\pi r^2 + \frac{128\pi}{r}}
        \]
    \end{itemize}

    \item \textbf{Dimensiones óptimas:}
    \begin{itemize}[leftmargin=1.5em, itemsep=2pt]
        \item Primera derivada:
        \[
        C'(r) = 16\pi r - \frac{128\pi}{r^2} = 0 \implies 16\pi r^3 = 128\pi \implies r^3 = 8 \implies \mathbf{r = 2\text{ m}}
        \]
        \item Altura correspondiente:
        \[
        h = \frac{16}{2^2} = \mathbf{4\text{ m}}
        \]
    \end{itemize}

    \item \textbf{Justificación y coste mínimo:}
    \begin{itemize}[leftmargin=1.5em, itemsep=2pt]
        \item Segunda derivada: $C''(r) = 16\pi + \frac{256\pi}{r^3} \implies C''(2) = 16\pi + \frac{256\pi}{8} = 48\pi > 0 \implies \textbf{Mínimo}$.
        \item Coste mínimo: $C(2) = 8\pi(2^2) + \frac{128\pi}{2} = 32\pi + 64\pi = \mathbf{96\pi\text{ €} \approx 301{,}59\text{ €}}$.
    \end{itemize}
\end{enumerate}

\vspace{0.8cm}

\noindent\textbf{2.} Se desea fabricar un depósito cilíndrico cerrado con una capacidad de $V = 54\pi\text{ m}^3$. Materiales: tapa superior a $4\text{ €/m}^2$, base inferior a $8\text{ €/m}^2$ y pared lateral a $6\text{ €/m}^2$.
\begin{enumerate}[label=\textbf{\alph*)}, leftmargin=2em, itemsep=6pt]
    \item \textbf{Función del coste total $C(r)$:}
    \begin{itemize}[leftmargin=1.5em, itemsep=2pt]
        \item Ligadura del volumen: $V = \pi r^2 h = 54\pi \implies h = \frac{54}{r^2}$.
        \item Coste total:
        \[
        C(r) = 4(\pi r^2) + 8(\pi r^2) + 6(2\pi r h) = 12\pi r^2 + 12\pi r \left(\frac{54}{r^2}\right) = \mathbf{12\pi r^2 + \frac{648\pi}{r}}
        \]
    \end{itemize}

    \item \textbf{Dimensiones óptimas:}
    \begin{itemize}[leftmargin=1.5em, itemsep=2pt]
        \item Primera derivada:
        \[
        C'(r) = 24\pi r - \frac{648\pi}{r^2} = 0 \implies 24\pi r^3 = 648\pi \implies r^3 = 27 \implies \mathbf{r = 3\text{ m}}
        \]
        \item Altura correspondiente:
        \[
        h = \frac{54}{3^2} = \mathbf{6\text{ m}}
        \]
    \end{itemize}

    \item \textbf{Justificación y coste mínimo:}
    \begin{itemize}[leftmargin=1.5em, itemsep=2pt]
        \item Segunda derivada: $C''(r) = 24\pi + \frac{1296\pi}{r^3} \implies C''(3) = 24\pi + \frac{1296\pi}{27} = 72\pi > 0 \implies \textbf{Mínimo}$.
        \item Coste mínimo: $C(3) = 12\pi(3^2) + \frac{648\pi}{3} = 108\pi + 216\pi = \mathbf{324\pi\text{ €} \approx 1017{,}88\text{ €}}$.
    \end{itemize}
\end{enumerate}

\end{document}
